\chapter{Conclusions and Future Work}
\label{chp:conclusions}

In this thesis, we ask a central question for autonomous drone deployment: how can autonomous navigation be provided on a weight-, size-, and power-constrained micro-drone with minimal changes to the environment? We address this question with a visible-light navigation system that uses lightweight on-board light sensing to guide the drone toward modulated light beacons, requiring only minimal modification to existing lighting infrastructure. To make the system robust in cluttered indoor environments, we extend it so the drone can visit several beacons in succession and reach a beacon whose direct light an obstacle partly blocks.

The work builds on the light-guided navigation system proposed by Huang et al.~\cite{huang2026lta}, especially their bearing-angle estimation and gradient-based light-source seeking algorithms. The bearing estimate provides an immediate direction toward a visible beacon, but it keeps no memory of the space the drone has already observed. The gradient method, in contrast, can follow changes in the light field, but it is slower because it must infer direction from measurements gathered during motion. These methods are insufficient for the more realistic missions considered in this thesis. They do not by themselves provide robust sequencing across multiple beacons. A bearing-only controller can also fail when an obstacle weakens or blocks the direct light signal, even though the bearing is the most efficient method in open-space source seeking.

This thesis set out to overcome these limitations within the constraints of a micro-drone. Meeting its two goals required solving four challenges, which follow from the three research questions of Chapter~\ref{chp:introduction}. Table~\ref{tab:rq-mapping} maps the research questions onto these challenges and goals. The following paragraphs discuss each challenge and how this thesis addresses it.

\begin{table}[htbp]
    \centering
    \caption{How the three research questions of Chapter~\ref{chp:introduction} map onto the four challenges solved in this thesis and the two navigation goals they enable. The first research question concerns feasibility and underpins both goals; the second and third are the goals themselves, each resting on one further challenge.}
    \label{tab:rq-mapping}
    \begin{tabular}{p{0.30\textwidth} p{0.34\textwidth} p{0.24\textwidth}}
        \hline
        Research question & Challenge & Goal \\
        \hline
        RQ1: Can the drone carry the light sensing on board and recover a reliable beacon signal in flight?
            & Challenge 1: Fit the light sensing within the drone's budget. \newline
              Challenge 2: Recover a reliable beacon signal in flight.
            & Foundation for RQ2 and RQ3 \\
        \hline
        RQ2: Can it navigate to several beacons in succession?
            & Challenge 3: Distinguish beacons by modulation frequency.
            & Goal 1: Visit several beacons in sequence \\
        \hline
        RQ3: Can it navigate around an obstacle blocking the beacon?
            & Challenge 4: Build a spatial memory of the light field.
            & Goal 2: Reach a beacon despite an obstacle \\
        \hline
    \end{tabular}
\end{table}

The first challenge is whether the sensing is feasible at all. We find that it is not the limiting factor. The custom board of Chapter~\ref{chp:hardware} provides eight-direction light sensing for under half of the drone's payload budget. It carries no processor of its own and adds no functionality the drone does not need. Useful directional light sensing is therefore well within the reach of a micro-drone. The difficulty lies in the navigation, not in the sensor.

The second challenge is to recover a trustworthy beacon signal in flight. Chapter~\ref{chp:realworld} develops the signal pipeline that does this. The decisive factor is how the pipeline estimates the noise floor. Using a median rather than a mean more than doubles the in-flight signal-to-noise ratio, from about 15 to about 32. The hardware anti-alias filter matters too, but only in combination. On its own it barely changes the ratio, yet removing it from the full pipeline lowers the in-flight ratio from about 32 to 22. Recovering the signal in flight therefore depends less on the measurement itself than on rejecting the noise around it.

The third challenge supports the first goal: distinguishing several beacons. We modulate each beacon at its own frequency, so a single recording carries every beacon as a separate peak, and the controller chooses its target by frequency. With this, the drone reaches several beacons in succession in real flight, as Chapter~\ref{chp:multibeacon} demonstrates. The first goal is therefore met.

The fourth challenge supports the second goal, which is the central problem of the thesis: reaching a beacon when an obstacle blocks the direct path. The enabling idea is spatial memory. The fused controller of Chapter~\ref{chp:algorithm} builds a map of the light field as it flies and estimates a gradient from it. The drone therefore retains a direction to follow even when the beacon is hidden. The contrast with the prior methods is sharp. As Chapter~\ref{chp:obstacle} shows, the bearing-only controller flies straight into the obstacle, and the single-photodiode gradient stalls against it. Both fail immediately. The fused controller instead detects when the beacon is occluded and curves around the obstacle without stopping, reaching the beacon in both simulation and real flight. The two real flights show this in two regimes. In the first, the bearing is lost behind the obstacle, and the drone navigates on the stored map alone. In the second, the blended heading alone is enough to round the obstacle, without ever falling back on the map. The second goal is therefore met: the drone keeps its efficiency while avoiding the obstacle, rather than crashing in a cluttered environment.

Underpinning both goals is the fused controller, the core technical contribution of this thesis. It adds a spatial memory of the light field to the bearing and gradient primitives of Huang et al.~\cite{huang2026lta}. At each moment it decides which of the two estimates to trust.

Taken together, the four challenges are solved, and the two goals with them, in both simulation and real flight. Huang et al.~\cite{huang2026lta} reached a single beacon in open space. This thesis reaches several in succession and routes around an obstacle that hides one, using only the lamps already in the room.

The result is not limited to one airframe. In a collaboration with Huang et al., whose blimp and simulator this work builds on, the same board and controller also navigate on their blimp. The navigation therefore depends on the light sensing, not on the platform.

Light-guided navigation thus moves a step closer to a practical, low-infrastructure positioning method for the smallest indoor drones.

\section{Future Work}
\label{sec:futurework}

Some limitations remain in the current system, and each points to a concrete direction for future work, so we treat the two together. We group them into improvements to the sensing hardware, extensions to the navigation method, and a broader characterisation and robustness study.

\subsection{Improving the Sensing Hardware}
\label{sec:fw-hardware}

The most immediate limitation is the system's operating range, which the photodiode sensing circuit bounds at both ends. When the drone is far from a beacon, the modulated signal becomes weak relative to steady background light. The SNR then falls below the level needed for reliable navigation, as measured in Chapter~\ref{chp:realworld}. When the drone is very close to the beacon, the photodiodes instead approach saturation and the recovered signal flattens. The dip seen near a strong beacon in Section~\ref{sec:snr-field} suggests this upper bound. The usable region therefore lies between these two limits. This limitation comes from the current circuit, not from visible-light navigation itself, and a redesigned front end could widen the range at both ends.

A transimpedance amplifier with adjustable or switched gain would let the front end follow the signal as the drone moves. Near a beacon, where the photodiodes now saturate, it would back off; far from it, where the signal is weak, it would amplify, widening the clean band at both ends. Higher-order band-pass filters tuned to the beacon band would help further, passing the modulation more selectively and rejecting the ambient light that otherwise reduces the sensor's usable range, the effect measured in Section~\ref{sec:ambient}.

The motor-noise measurements of Chapter~\ref{chp:obstacle}, reported in Section~\ref{sec:motor-noise}, show the CrazyFlie's motors are the main cost of flight, and that vibration and electrical interference contribute in roughly equal measure. This opens a clear hardware direction. Soft-mounting the board or the motors to damp the vibration should recover the mechanical half of the loss, while better shielding and supply filtering would address the electrical half. The same measurements explain why the blimp records a higher SNR than the CrazyFlie: as Section~\ref{sec:blimp-generalisation} shows, the blimp runs its motors far less. A controlled side-by-side flight of the two carrying the same board would then measure how much of that gap soft-mounting and shielding can close.

Photodiode geometry is a further direction, and it is also what restricts the system to fixed-altitude flight. The photodiodes sit in a horizontal ring, as shown in Section~\ref{sec:pcb-layout}. The bearing therefore only resolves direction within that plane, and the array responds weakly to light from above or below it. The ring was chosen deliberately, to reduce the overhead ambient light entering the sensors. A more spherical arrangement would let the system sense the vertical directions too. It would admit more ambient light in exchange, which a more capable sensing circuit could manage.

Covering elevation carries a hardware cost. A single ring uses $N$ photodiodes, with $N=8$ in Chapter~\ref{chp:hardware}. Covering the upper hemisphere at the same angular spacing would need several such rings stacked over elevation, so the sensor count grows roughly with the square of the per-ring count. Full spherical coverage is not required in practice, since the drone need not sense light from directly below. This keeps the count well under that worst case.

The same redesign would also support faster, more agile flight. During aggressive manoeuvres the drone tilts, changing where each photodiode points. In this thesis we keep the flight gentle and near-level, so that the horizontal-plane bearing of Section~\ref{sec:bearing} stays valid. Future work could instead measure the tilt and correct the bearing error that more aggressive motion introduces.

\subsection{Extending the Navigation}
\label{sec:fw-navigation}

The controller currently acts only in the horizontal plane, since the algorithms of Chapter~\ref{chp:algorithm} use a two-dimensional light map indexed by the drone's horizontal position. Lifting the navigation into three dimensions would let the drone seek beacons above or below its flight height. The algorithmic extension is relatively direct. The light map of Section~\ref{sec:light-intensity-map} would store intensity over a volume, indexed by horizontal position and altitude rather than over a flat grid. This raises it from an $O(G^2)$ grid to an $O(G^3)$ volume for $G$ cells per axis. The gradient estimator of Section~\ref{sec:gradient-estimation} would simply add a third coordinate. This increases computation, but not prohibitively. The harder part is sensing: the extension needs light variation along the elevation axis, which the horizontal photodiode ring cannot provide. For this reason, three-dimensional navigation should be developed together with the photodiode-board redesign above, since elevation sensing cannot come from software alone.

Future work should also test the controller in more diverse obstacle settings than the single box used in Chapter~\ref{chp:obstacle}. Real deployments may contain pillars, tables, shelves, walls, or other objects that block and reshape the light field in different ways. Multiple obstacles, narrow corridors, and concave layouts could each stress the map-based recovery mechanism of Section~\ref{sec:state-machine} by creating shadows, detours, or local maxima that the stored gradient may not escape. Evaluating these cases would clarify when the current recovery strategy is sufficient and when additional planning or obstacle awareness is needed.

A further direction is to relax the static-environment assumption. All experiments in this thesis keep the beacons and obstacles fixed while the drone navigates. The map is rebuilt from scratch for each beacon, as Section~\ref{sec:light-intensity-map} describes. A change between one beacon and the next is therefore absorbed when the map is reset. The harder case is a beacon or obstacle that moves \emph{during} a single approach, once the drone has already mapped part of the light field. The affected cells then hold stale values, and the controller may follow an outdated gradient until it remaps the region. Handling this robustly would let the controller detect that a stored cell no longer matches the field and refresh it on the fly, rather than relying on the per-beacon reset alone. Evaluating the controller under such dynamic conditions would clarify when the approach remains effective.

Another direction is to refine how the controller fuses the bearing and gradient estimates. The current controller uses fixed, equal weights whenever both estimates are valid, as defined by the fusion rule in Equation~\ref{eq:fusion}. When only one estimate remains valid, the state machine mentioned in Section~\ref{sec:state-machine} switches to only using that estimate. Our design is simple, but it treats each input as either trusted or ignored, even though both estimates already carry information about their reliability. A more flexible approach would be for the controller to adjust the bearing and gradient weights, $w_B$ and $w_G$, according to their confidence. The bearing confidence could come from the summed photodiode magnitude used as the validity score in Section~\ref{sec:bearing-angle}, while the gradient confidence could come from the $R^2$ value of the plane fit in Equation~\ref{eq:gradient-r2}. With this design, the drone would rely more on the bearing when the beacon is clearly visible and more on the gradient when the bearing weakens inside an obstacle shadow, while still allowing both estimates to contribute to the final heading. Such confidence-weighted fusion would make the transition between the two estimators smoother than the current hard-threshold hand-off.

\subsection{Characterisation and Robustness}
\label{sec:fw-characterisation}

The real-world results in this thesis are demonstrations rather than a measure of reliability. We flew more missions than the few shown in Chapters~\ref{chp:multibeacon} and~\ref{chp:obstacle}, and the drone reached its beacon on most of them, but we did not run a controlled repeatability study. This matters for deployment: a system trusted to navigate on its own must be characterised by how often it succeeds, not only by whether it can. Quantifying that success rate across varied conditions is the broad goal of the study below.

The controller contains several tunable parameters, listed in Appendix~\ref{app:parameters}. Their values are currently selected empirically for the experiments in this thesis, so future work should quantify how sensitive navigation performance is to each setting. A formal ablation study would address this by repeating the same missions under fixed conditions while varying one parameter at a time and keeping the others unchanged. This study would also provide the controlled evaluation the present demonstrations lack. By measuring success rate, path length, travel time, and recovery behaviour under each parameter setting, such an ablation study would replace the current hand-tuned values with practical tuning guidelines for deploying the system in new environments.

The arrival condition deserves closer attention before real deployment. In the current state machine, described in Section~\ref{sec:state-machine}, the controller registers arrival using a single SNR threshold. The three-waypoint flight in Section~\ref{sec:multibeacon-traverse} shows the weakness of this design: although the SNR rises above the threshold near the final beacon, the state machine does not reliably latch the arrival event, and the drone continues flying. A more robust arrival mechanism should therefore avoid relying on an instantaneous SNR value alone. For example, the controller could require the SNR to stay above the threshold for several consecutive frames before confirming arrival, or estimate the drone's position relative to the target beacon. This would make beacon handover more reliable in longer missions.

